\documentclass[11pt,a4paper]{article}

\usepackage[margin=2.5cm]{geometry}
\usepackage{amsmath,amssymb,amsfonts}
\usepackage{booktabs}
\usepackage{enumitem}
\usepackage{mathtools}
\usepackage{xcolor}
\usepackage{hyperref}

\title{Formal ILP Formulation for RunSoC 2.0}

\begin{document}
\maketitle

\section{Problem Definition}

This document formalizes the integer linear programming (ILP) model used for assigning and scheduling a set of tasks and their associated execution instances (jobs) on a heterogeneous multi-core platform.
The model assumes a \emph{partitioned scheduling} approach, where all jobs belonging to a specific task must execute on the same core.
The model minimizes a weighted objective consisting of memory-budget overflow penalties and communication penalties between dependent tasks, subject to strict precedence and timing constraints.

\section{Sets}

\begin{align*}
T &:= \text{set of tasks},\\
J &:= \text{set of jobs (task instances)},\\
C &:= \text{set of cores},\\
K &:= \text{set of clusters},\\
D_{\mathrm{job}} &\subseteq J \times J := \text{set of job execution precedence dependencies},\\
D_{\mathrm{task}} &\subseteq T \times T := \text{set of task dependencies (for evaluating communication penalties)},\\
E_t &\subseteq C := \text{set of eligible cores for task } t \in T,\\
C_k &:= \{c \in C \mid c \text{ belongs to cluster } k\}, \quad \forall k \in K.
\end{align*}

Let $\tau(j) \in T$ denote the parent task to which job $j \in J$ belongs.
A dependency $(j, j') \in D_{\mathrm{job}}$ means that job $j$ must finish before job $j'$ may start.
A dependency $(t_1, t_2) \in D_{\mathrm{task}}$ indicates data flow between task $t_1$ and task $t_2$, incurring a penalty if assigned to different cores or clusters.

\section{Parameters}

For each job $j \in J$:
\begin{align*}
r_j &:= \text{release time of job } j,\\
d_j &:= \text{absolute deadline of job } j \text{ (if specified)},\\
p_j &:= \text{base processing duration of job } j.
\end{align*}

For each task $t \in T$:
\begin{align*}
m_t &:= \text{memory demand of task } t.
\end{align*}

For each core $c \in C$:
\begin{align*}
\alpha_c &:= \text{WCET scaling factor of core } c,\\
B_c^{\mathrm{core}} &:= \text{memory budget of core } c.
\end{align*}

For each cluster $k \in K$:
\begin{align*}
B_k^{\mathrm{cluster}} &:= \text{memory budget of cluster } k.
\end{align*}

Let $\kappa(c) \in K$ denote the cluster containing core $c$.

The processing time of job $j$ on core $c$ is:
\begin{equation}
p_{jc} := p_j \alpha_c.
\end{equation}

The model utilizes a sufficiently large constant $M$ for linearizing non-overlap constraints:
\begin{equation}
M = \max_{j \in J} d_j \quad (\text{or maximal achievable horizon}).
\end{equation}

Penalty weights are denoted by:
\begin{align*}
\lambda_{\mathrm{core}} &:= \text{penalty weight for core-memory overflow},\\
\lambda_{\mathrm{cluster}} &:= \text{penalty weight for cluster-memory overflow},\\
w_{\mathrm{intra}} &:= \text{penalty for communication on the same core},\\
w_{\mathrm{intercore}} &:= \text{penalty for communication between different cores in the same cluster},\\
w_{\mathrm{intercluster}} &:= \text{penalty for communication between cores in different clusters}.
\end{align*}

For an explicitly defined communication path $(c,d)$, let $q_{cd}^{\mathrm{explicit}}$ be its specified penalty. The effective communication penalty $q_{cd}$ between cores $c,d \in C$ is:
\begin{equation}
q_{cd} =
\begin{cases}
q_{cd}^{\mathrm{explicit}}, & \text{if an explicit communication path } (c,d) \text{ is defined},\\
w_{\mathrm{intra}}, & \text{if } c=d,\\
w_{\mathrm{intercore}}, & \text{if } c \neq d \text{ and } \kappa(c)=\kappa(d),\\
w_{\mathrm{intercluster}}, & \text{if } \kappa(c)\neq \kappa(d).
\end{cases}
\end{equation}

\section{Decision Variables}

\begin{align*}
y_{tc} &\in \{0,1\}
&& \forall t \in T,\ c \in C,
&& \text{equals 1 iff task } t \text{ is assigned to core } c,\\
x_{jc} &\in \{0,1\}
&& \forall j \in J,\ c \in C,
&& \text{equals 1 iff job } j \text{ is assigned to core } c,\\
s_j &\geq 0
&& \forall j \in J,
&& \text{start time of job } j,\\
f_j &\geq 0
&& \forall j \in J,
&& \text{finish time of job } j,\\
o_c^{\mathrm{core}} &\geq 0
&& \forall c \in C,
&& \text{memory overflow on core } c,\\
o_k^{\mathrm{cluster}} &\geq 0
&& \forall k \in K,
&& \text{memory overflow on cluster } k.
\end{align*}

For each unordered pair of distinct jobs $\{j, j'\}$ with $j < j'$:
\begin{align*}
u_{jj'} &\in \{0,1\}
&& \forall j, j' \in J,\ j<j',
&& \text{ordering variable used for non-overlap constraints}.
\end{align*}

For each task dependency $(t_1,t_2)\in D_{\mathrm{task}}$ and eligible core pair $c \in E_{t_1}$, $d \in E_{t_2}$:
\begin{align*}
z_{t_1 t_2 cd} &\in \{0,1\}
&& \forall (t_1,t_2)\in D_{\mathrm{task}},\ c\in E_{t_1},\ d\in E_{t_2},
&& \text{equals 1 iff } t_1 \text{ is assigned to } c \text{ and } t_2 \text{ is assigned to } d.
\end{align*}

\section{Objective Function}

The objective minimizes weighted core-memory overflow, weighted cluster-memory overflow, and communication penalties along task dependency edges. Schedule makespan is managed strictly through release times and deadlines rather than objective optimization.

\begin{equation}
\min
\quad
\lambda_{\mathrm{core}}
\sum_{c \in C} o_c^{\mathrm{core}}
+
\lambda_{\mathrm{cluster}}
\sum_{k \in K} o_k^{\mathrm{cluster}}
+
\sum_{(t_1,t_2)\in D_{\mathrm{task}}}
\sum_{c \in E_{t_1}}
\sum_{d \in E_{t_2}}
q_{cd} z_{t_1 t_2 cd}.
\end{equation}

\section{Constraints}

\subsection{Task Assignment \& Partitioned Scheduling}

Each task must be assigned exactly once to one of its eligible cores:
\begin{equation}
\sum_{c \in E_t} y_{tc} = 1
\qquad
\forall t \in T.
\end{equation}

Assignments to ineligible cores are forbidden:
\begin{equation}
y_{tc}=0
\qquad
\forall t \in T,\ \forall c \in C \setminus E_t.
\end{equation}

To enforce partitioned scheduling, every job must execute on the core to which its parent task is assigned:
\begin{equation}
x_{jc} = y_{\tau(j) c}
\qquad
\forall j \in J,\ \forall c \in C.
\end{equation}

\subsection{Timing \& Deadline Constraints}

Each job may not start before its release time:
\begin{equation}
s_j \geq r_j
\qquad
\forall j \in J.
\end{equation}

If a chain root job has strict jitter constraints, its start time may optionally be bounded:
\begin{equation}
s_j \leq r_j + \text{jitter}
\qquad \text{(for applicable root jobs)}.
\end{equation}

The finish time of each job is determined by its start time and the processing time induced by the selected core:
\begin{equation}
f_j
=
s_j
+
\sum_{c \in C} p_j \alpha_c x_{jc}
\qquad
\forall j \in J.
\end{equation}

Terminal jobs in a chain with an defined absolute deadline must finish prior to that deadline:
\begin{equation}
f_j \leq d_j
\qquad
\forall j \in J \text{ with } d_j \neq \emptyset.
\end{equation}

\subsection{Precedence Constraints}

For every job execution dependency, the successor job may only start after the predecessor job has finished:
\begin{equation}
s_{j'} \geq f_j
\qquad
\forall (j,j') \in D_{\mathrm{job}}.
\end{equation}

\subsection{Memory-Budget Constraints}

Memory footprints are evaluated at the \emph{Task} level.

Core-level memory use may exceed the core memory budget only through the non-negative overflow variable:
\begin{equation}
\sum_{t \in T} m_t y_{tc}
\leq
B_c^{\mathrm{core}} + o_c^{\mathrm{core}}
\qquad
\forall c \in C.
\end{equation}

Cluster-level memory use may exceed the cluster memory budget only through the non-negative overflow variable:
\begin{equation}
\sum_{t \in T}
\sum_{c \in C_k}
m_t y_{tc}
\leq
B_k^{\mathrm{cluster}} + o_k^{\mathrm{cluster}}
\qquad
\forall k \in K.
\end{equation}

\subsection{Non-Overlap Constraints on Cores}

For every unordered pair of jobs $j,j' \in J$ with $j<j'$, and for every core $c \in C$, the model enforces that if both jobs are assigned to the same core, one must finish before the other starts.

\begin{equation}
s_{j'}
\geq
f_j
-
M\left(3 - u_{jj'} - x_{jc} - x_{j'c}\right)
\qquad
\forall j,j' \in J,\ j<j',\ \forall c \in C.
\end{equation}

\begin{equation}
s_j
\geq
f_{j'}
-
M\left(2 + u_{jj'} - x_{jc} - x_{j'c}\right)
\qquad
\forall j,j' \in J,\ j<j',\ \forall c \in C.
\end{equation}

Here, $u_{jj'}=1$ selects the ordering $j$ before $j'$, while $u_{jj'}=0$ selects the ordering $j'$ before $j$.

\subsection{Communication-Penalty Linearization}

For each task dependency $(t_1,t_2)\in D_{\mathrm{task}}$ and each eligible core pair $c\in E_{t_1}$, $d\in E_{t_2}$, the auxiliary variable $z_{t_1 t_2 cd}$ represents the logical conjunction $z_{t_1 t_2 cd} = y_{t_1 c} y_{t_2 d}$.

Because the objective minimizes the sum over $z_{t_1 t_2 cd}$ and all coefficients $q_{cd}$ are non-negative, the product can be sufficiently linearized with the single lower-bound constraint:
\begin{equation}
z_{t_1 t_2 cd} \geq y_{t_1 c} + y_{t_2 d} - 1
\qquad
\forall (t_1,t_2)\in D_{\mathrm{task}},\ c\in E_{t_1},\ d\in E_{t_2}.
\end{equation}

\section{Variable Domains}

\begin{align*}
y_{tc} &\in \{0,1\}
&& \forall t \in T,\ c \in C,\\
x_{jc} &\in \{0,1\}
&& \forall j \in J,\ c \in C,\\
u_{jj'} &\in \{0,1\}
&& \forall j, j' \in J,\ j<j',\\
z_{t_1 t_2 cd} &\in \{0,1\}
&& \forall (t_1,t_2)\in D_{\mathrm{task}},\ c\in E_{t_1},\ d\in E_{t_2},\\
s_j, f_j &\in [0,M]
&& \forall j \in J,\\
o_c^{\mathrm{core}} &\geq 0
&& \forall c \in C,\\
o_k^{\mathrm{cluster}} &\geq 0
&& \forall k \in K.
\end{align*}

\section{Remarks}

\begin{itemize}
    \item \textbf{Decoupling of Tasks and Jobs:} The current architecture decouples memory footprints (Tasks) from temporal executions (Jobs). While multiple instances (Jobs) of a Task may exist in the schedule, the core's memory budget is only impacted once per assigned Task.
    \item \textbf{Constraint Solvers:} In practice, tools like CP-SAT natively handle intervals and `NoOverlap` constraints, removing the need for manual Big-M transformations ($u_{jj'}$). This document represents the formal mathematical equivalent.
    \item \textbf{Memory-Pressure Optimization Weight:} Memory budgets remain modeled as soft constraints to preserve model feasibility. The scalar weights $\lambda_{\mathrm{core}}$ and $\lambda_{\mathrm{cluster}}$ define the preferred optimization focus on memory pressure: values closer to $0$ reduce the importance of memory-budget violations in the objective, while larger values increase the optimizer's preference for assignments that respect core- and cluster-level memory limits.
\end{itemize}

\end{document}